\chapter{Hand and Finger Fractures}

\section*{Definition and Overview}
Hand and finger fractures are breaks in the bones of the hand, including the metacarpals (the long bones of the palm) and the phalanges (the finger bones). They are among the most common fractures, often caused by falls onto an outstretched hand, direct blows, sports injuries, and punching a hard object. Symptoms include pain, swelling, bruising, deformity, and difficulty moving the finger or hand. Most hand and finger fractures are treated without surgery, with splinting or buddy strapping, and heal within four to six weeks. Some fractures, particularly those that are displaced, unstable, or involve joints, need surgical fixation. Hand function is usually fully restored with treatment and hand therapy.

\section*{Epidemiology}
Hand and finger fractures are extremely common, accounting for around 10--20% of all fractures presenting to emergency departments. They are most frequent in young adult men (often from sport or punching), in children (falls), and in older people (falls with osteoporosis). The fifth metacarpal ("boxer's fracture") and the distal phalanx are among the most commonly broken bones. Most hand fractures are simple and heal well, but up to a third involve joints or are displaced enough to need surgery.

\section*{Aetiology and Pathophysiology}
Hand fractures result from direct or indirect force to the bones of the hand.

\begin{itemize}
    \item \textbf{Mechanism:} falls onto an outstretched hand, direct blows, crushing injuries, and punching
    \item \textbf{Common fractures:} distal phalanx (fingertip), phalangeal shaft, metacarpal neck (boxer's fracture), and intra-articular fractures involving joints
    \item \textbf{Displacement:} broken ends may shift out of alignment
    \item \textbf{Angulation and rotation:} fractures can angulate or rotate, causing finger overlap or scissoring
    \item \textbf{Open fractures:} when the bone pierces the skin, requiring urgent treatment
    \item \textbf{Associated injuries:} tendon and ligament injuries, nail bed injuries, and soft tissue damage
    \item \textbf{Bone healing:} most hand bones heal well because of excellent blood supply, but stiffness is the main challenge
\end{itemize}

\section*{Clinical Features}
The features of a hand or finger fracture are usually obvious.

\begin{itemize}
    \item Pain at the site of the injury, worse with movement or touch
    \item Swelling and bruising
    \item Tenderness over the bone
    \item Deformity, such as a shortened, angulated, or rotated finger
    \item Difficulty moving the finger or gripping
    \item Numbness or tingling if nerves are affected
    \item A wound near the fracture in open fractures
    \item The finger may be out of alignment with the others
\end{itemize}

\section*{Differential Diagnosis}
Other hand injuries can mimic fractures.

\begin{itemize}
    \item Sprains and ligament injuries
    \item Tendon injuries, including mallet finger (ruptured extensor tendon) and jersey finger (ruptured flexor tendon)
    \item Dislocations of the finger joints
    \item Contusions (bruising) of the hand
    \item Nail bed injuries
    \item Fractures with tendon involvement, which may be missed
    \item In children, growth plate injuries
\end{itemize}

\section*{When to Seek Medical Care}
Seek medical care for any suspected hand or finger fracture, especially if there is deformity, numbness, an open wound, or an inability to move the finger. Early treatment prevents stiffness and deformity.

\textbf{Call 000 (or local emergency services) for:}
\begin{itemize}
    \item A finger that is pale, cold, blue, or numb, which may indicate a vascular injury
    \item An open fracture with bone visible or a deep wound
    \item A fractured or amputated part, which may be replantable
    \item Severe bleeding that does not stop with pressure
    \item Signs of severe infection, including spreading redness and fever
\end{itemize}

\section*{Diagnosis and Investigations}
Diagnosis is clinical and confirmed with X-rays.

\begin{itemize}
    \item \textbf{Examination:} inspection, palpation, and assessment of movement, sensation, and circulation
    \item \textbf{X-rays:} confirm the fracture, its location, displacement, and any joint involvement
    \item \textbf{Multiple views:} oblique views help detect subtle fractures
    \item \textbf{Assessment of rotation:} checking finger alignment when the hand is in a fist
    \item \textbf{CT:} rarely needed, but useful for complex intra-articular fractures
    \item \textbf{Follow-up X-rays:} monitor healing
\end{itemize}

\section*{Management}
Most hand and finger fractures are managed without surgery.

\begin{itemize}
    \item \textbf{Splinting:} a splint or cast immobilises the fracture and allows healing
    \item \textbf{Buddy strapping:} taping the injured finger to its neighbour is used for many stable phalangeal fractures
    \item \textbf{Reduction:} displaced fractures may be manipulated back into alignment under local anaesthetic
    \item \textbf{Surgery:} pins, plates, or screws for displaced, unstable, or joint fractures
    \item \textbf{Pain relief:} paracetamol or ibuprofen
    \item \textbf{Elevation and ice:} reduce swelling in the first days
    \item \textbf{Hand therapy:} exercises to prevent stiffness once the fracture is stable
    \item \textbf{Follow-up:} repeat X-rays and review to confirm healing
    \item \textbf{Return to activity:} gradual return to sport and heavy use after healing
\end{itemize}

\section*{Complications}
\begin{itemize}
    \item Stiffness, the most common complication, especially in older people
    \item Malunion, in which the bone heals in a deformed position
    \item Non-union, in which the bone fails to heal
    \item Rotational deformity causing finger scissoring
    \item Post-traumatic arthritis, especially with joint fractures
    \item Tendon adhesions and reduced movement
    \item Infection in open fractures or after surgery
    \item Nerve injury causing numbness or weakness
    \item Complex regional pain syndrome, a rare complication with persistent pain and swelling
\end{itemize}

\section*{Prognosis and Outcomes}
Hand and finger fractures generally heal well, and most people regain full or near-full function. Stable fractures treated with splinting or buddy strapping heal within four to six weeks, and hand therapy restores movement. Surgical fixation produces good results for the fractures that need it. Stiffness is the main concern and is prevented by early, guided movement once the bone is stable. Most people return to work, sport, and daily activities within two to three months.

\section*{Prevention and Self-Care}
\begin{itemize}
    \item Wear appropriate protective equipment for sport, including gloves and wrist guards
    \item Use proper technique when punching or striking, and avoid punching hard surfaces
    \item Keep floors and walkways free of trip hazards
    \item Maintain bone health with calcium, vitamin D, and weight-bearing exercise
    \item Seek early treatment for hand injuries
    \item Follow splinting and activity instructions
    \item Do hand exercises as directed by your therapist
    \item Elevate the hand and use ice in the first days
    \item Attend follow-up appointments and X-rays
\end{itemize}

\section*{Sources and Further Reading}
The following reputable sources provide further information for healthcare students and patients seeking reliable, current advice about hand and finger fractures.

\begin{itemize}
    \item Australian Orthopaedic Association. \textit{Fracture information for patients.} https://aoa.org.au/
    \item Australian Hand Surgery Society. \textit{Hand fractures and treatment.} https://handsurgerysociety.org.au/
    \item Healthdirect Australia. \textit{Hand and finger fractures.} https://www.healthdirect.gov.au/
    \item Sports Medicine Australia. \textit{Hand injuries fact sheet.} https://sma.org.au/
    \item NHS. \textit{Hand and finger fractures.} https://www.nhs.uk/
    \item American Academy of Orthopaedic Surgeons. \textit{Hand fractures.} https://orthoinfo.aaos.org/
\end{itemize}
